% Secao 7 - Validacao externa contra a regionalizacao de saude (CIR/regiao de
% saude). Disciplina: descritivo. Numeros com % src: 14_external_validation.py

\section{Validação externa: o fluxo e a regionalização do SUS}
\label{sec:validacao}

Uma medida de deserto construída a partir do próprio fluxo corre o risco de ser
circular. Para testá-la contra um objeto independente, confrontamos os dois
conceitos com a regionalização do SUS --- as \valNregioes{} regiões de saúde que agrupam
municípios e que existem, por desenho, para conter a referência: o paciente
deveria ser atendido dentro da própria região. O fluxo revelado permite medir
quanto essa promessa se cumpre.

O primeiro resultado é que ela se cumpre pela metade. Apenas \valWithinMean\% das
internações ocorrem dentro da região de saúde do município de residência; a
regionalização administrativa contém pouco mais da metade do fluxo real.% src: 14_external_validation - mean_within 0.564
E a parte que escapa não é uniforme. Nos municípios perto-mas-isolados, apenas
\valWithinFO\% das internações ficam na própria região, contra \valWithinNE\% nos
municípios conectados, e \valTophubOutFO\% deles têm seu principal hub externo
\emph{fora} da região de saúde designada.% src: by_quadrant - pct_within_region 46.2 vs 71.2; pct_tophub_out 65.4
Os desertos de fluxo são, precisamente, os municípios em que o desenho de
regionalização falha em conter a referência.

O segundo resultado é qual conceito recupera melhor essa falha. O burden de fluxo
F1 correlaciona-se com o escape da região ($r=\valCorrWithinF$ entre contenção
intra-regional e F1) com mais força do que a distância ao hospital mais próximo
($r=\valCorrWithinIso$).% src: corr - r_within_f1 -0.274, r_within_iso -0.126
Ou seja, a medida de fluxo não apenas é independente da regionalização ---
construída sem qualquer referência a ela --- como reconstrói, melhor que a
distância, onde essa regionalização não se sustenta na prática. É a validação que
uma medida derivada do próprio fluxo precisa: ela recupera um objeto
administrativo que não entrou em sua construção.

Uma ressalva, na mesma chave das anteriores, disciplina a leitura. Referência
para fora da região nem sempre é falha: parte da alta complexidade é,
legitimamente, concentrada em centros que servem várias regiões. O baixo
\emph{within-region} de um deserto de fluxo sinaliza um descompasso entre o mapa
administrativo e o uso revelado a investigar, não um veredito automático de má
alocação. O que a Seção estabelece é que F1 enxerga esse descompasso, e o enxerga
melhor que a régua de quilômetros.
